% !TeX TXS-program:compile = txs:///pdflatex/[-shell-escape]
% !TeX TXS-program:pdflatex = pdflatex -synctex=1 -interaction=nonstopmode --shell-escape %.tex
\section{Control Flow}\label{sec:control_flow}

\subsection{Video - Registers}\label{subsec:video_registers}

\subsection{Video - Control Flow}\label{subsec:video_control_flow}

\subsection{Comparing Values} \label{subsec:comparing_values}

Para poder realizar las comparaciones, se utiliza la instrucción \textbf{\textit{cmp}},
este compara 2 valores restando el segundo operando del primero.
El valor de dicha resta, no se almacena en ningún lugar al que se pueda
acceder directamente, este actualiza las \textbf{\textit{flag}} internas de la CPU.

\begin{codebox}[ASM]{Ejemplo}
	cmp rdi, 42             ; Se comprar un valor inmediato
	cmp QWORD PTR [rsp], 42 ; tambien se puede comparar direcciones de memoria
\end{codebox}

Esto de aquí, calcula internamente \( rdi - 42 \) y la CPU estable un bit
especial llamado \textbf{\textit{Zero Flag(ZF)}}, en caso de que la resta
sea \( 0 \) el valor de \textbf{\textit{ZF}} es \( 1 \), caso contrario va a tomar el valor de \( 0 \).

Con lo anteriormente mencionado x86 nos proporciona una familia de instrucciones
que escriben \( 0 \) o \( 1 \) en un destino de 1 byte.

\begin{codebox}[ASM]{}
	setz dil
\end{codebox}

Esto comprueba el valor de  \textbf{\textit{ZF}}, y lo escribe en \textbf{\textit{dil}}.
Y a su vez, \textbf{\textit{setnz(Set if Not Zero)}} que hace lo contrario.

Recordar que \textbf{\textit{dil}} son los primeros \( 8 \) bits menos significativos
del registro \textbf{\textit{rdi}}.
\begin{codebox}[ASM]{Solución del Problema}
	.intel_syntax noprefix
	cmp QWORD PTR [rsp], 42
	setz dil
	mov rax, 60
	syscall
\end{codebox}
\begin{flagbox}{Flag Capturada}
	pwn.college{s9M8ey0NvS8aym8WgNaOwwTzosR.01M0czMywSMycjN4EzW}
\end{flagbox}

\subsection{Comparing Characters} \label{subsec:comparing_characters}

En el módulo de Stack, se vio que \mintinline{ASM}{[rsp + 16]}
es un puntero al primero valor
pasado como parámetro \textbf{\textit{argv[1]}}. Y como vimos en el capitulo anterior \ref{sec:output_and_input}
podemos solo utilizar el valor de un solo byte (\mintinline{ASM}{BYTE PTR [rax]}), y con esto podemos realizar
comparaciones de un solo caracter.

\begin{codebox}[ASM]{Solución}
	.intel_syntax noprefix
	mov rax, [rsp + 16]
	cmp BYTE PTR [rax], 'p'
	setz dil
	mov rax, 60
	syscall
\end{codebox}

\begin{flagbox}{Flag Capturada}
	pwn.college{EK0Xt5rxYiy5yQWdKCGSVY3dSbL.0FN0czMywSMycjN4EzW}
\end{flagbox}

\subsection{Conditional Control Flow } \label{subsec:conditional_control_flow}

Con el resultado anterior, se le puede indicar a la CPU que salte a otra parte
del código.

Estos se ejecutan después de la instrucción \mintinline{ASM}{cmp}, si quiere
revisar todos los \href{https://www.felixcloutier.com/x86/jcc}{saltos condicionales}.

Aquí una tabla resumen:

\begin{table}[htpb]
	\centering
	\renewcommand{\arraystretch}{1.2}
	\begin{tabular}{|l|l|l|l|l|}
		\hline
		\textbf{Instrucción}        & \textbf{Mnemónico / Significado} & \textbf{Condición lógica}      & \textbf{Banderas evaluadas}             \\ \hline
		\texttt{jmp}                & Jump                             & Incondicional                  & Ninguna (salta siempre)                 \\ \hline
		\texttt{je} / \texttt{jz}   & Jump if Equal / Zero             & $A == B$ o resultado cero      & \texttt{ZF = 1}                         \\ \hline
		\texttt{jne} / \texttt{jnz} & Jump if Not Equal / Not Zero     & $A \neq B$ o resultado no cero & \texttt{ZF = 0}                         \\ \hline
		\texttt{jg} / \texttt{jnle} & Jump if Greater                  & $A > B$ (con signo)            & \texttt{ZF = 0} y \texttt{SF = OF}      \\ \hline
		\texttt{jge} / \texttt{jnl} & Jump if Greater or Equal         & $A \ge B$ (con signo)          & \texttt{SF = OF}                        \\ \hline
		\texttt{jl} / \texttt{jnge} & Jump if Less                     & $A < B$ (con signo)            & \texttt{SF $\neq$ OF}                   \\ \hline
		\texttt{jle} / \texttt{jng} & Jump if Less or Equal            & $A \le B$ (con signo)          & \texttt{ZF = 1} o \texttt{SF $\neq$ OF} \\ \hline
		\texttt{ja} / \texttt{jnbe} & Jump if Above                    & $A > B$ (sin signo)            & \texttt{CF = 0} y \texttt{ZF = 0}       \\ \hline
		\texttt{jae} / \texttt{jnb} & Jump if Above or Equal           & $A \ge B$ (sin signo)          & \texttt{CF = 0}                         \\ \hline
		\texttt{jb} / \texttt{jc}   & Jump if Below / Carry            & $A < B$ (sin signo) / Acarreo  & \texttt{CF = 1}                         \\ \hline
		\texttt{jbe} / \texttt{jna} & Jump if Below or Equal           & $A \le B$ (sin signo)          & \texttt{CF = 1} o \texttt{ZF = 1}       \\ \hline
		\texttt{js}                 & Jump if Sign                     & Resultado negativo             & \texttt{SF = 1}                         \\ \hline
		\texttt{jns}                & Jump if Not Sign                 & Resultado positivo o cero      & \texttt{SF = 0}                         \\ \hline
		\texttt{jo}                 & Jump if Overflow                 & Desbordamiento con signo       & \texttt{OF = 1}                         \\ \hline
		\texttt{jno}                & Jump if Not Overflow             & Sin desbord. con signo         & \texttt{OF = 0}                         \\ \hline
	\end{tabular}
	\caption{ Algunos saltos en ensamblador x86-64}
	\label{tab:saltos_asm}
\end{table}

\begin{codebox}[ASM]{}
	cmp BYTE PTR [rax], 'p'
	jne fail
\end{codebox}

\textbf{\textit{fail}} es una \textbf{label}, estas no son una instrucción de máquina, sino como indicador para que el CPU
sepa a que parte tiene que saltar para continuar con la ejecución.

\begin{codebox}[ASM]{}
	fail:
	mov rdi, 1
	mov rax, 60
	syscall
\end{codebox}

En caso de que no se origine el salto, va a continuar con la ejecución normal.

\begin{codebox}[ASM]{Ejemplo de uso}
	main:
	[load and compare]
	jne fail          ←;jump to fail if NOT equal

	success:
	mov rdi, 0
	mov rax, 60
	syscall

	fail:
	mov rdi, 1
	mov rax, 60
	syscall
\end{codebox}

\begin{codebox}[ASM]{Solución del problema}
	.intel_syntax noprefix
	mov rax, [rsp + 16]
	cmp BYTE PTR [rax], 'p'
	jne fail

	mov rdi,0
	mov rax, 60
	syscall

	fail:
	mov rdi,1
	mov rax,60
	syscall
\end{codebox}

\begin{flagbox}{Flag Capturada}
	pwn.college{QZ9khdqQojrEPg_83_78vJ54Ecc.0VN0czMywSMycjN4EzW}
\end{flagbox}


\subsection{Comparing  Strings} \label{subsec:comparing_strings}

Para poder comparar string, se tiene que comparar letra por letra:
\begin{codebox}[ASM]{}
	mov rax, [rsp + 16]

	cmp BYTE PTR [rax], 'Y'
	jne fail

	cmp BYTE PTR [rax+1], 'E'
	jne fail

	cmp BYTE PTR [rax+2], 'S'
	jne fail
\end{codebox}
\begin{flagbox}{Flag Capturada}
	pwn.college{w_w8UjT_9EaU1ildN9bT0-Y5gF9.0lN0czMywSMycjN4EzW}
\end{flagbox}

\subsection{Reverse the Password} \label{subsec:reverse_the_password}

Este desafio busca recordar los temas anteriores, como el \mintinline{bash}{objdump}, y lo que hemos visto
anteriormennte para la comparación de caracteres y el poder leer código en \textbf{\textit{ASM}}.

\begin{infobox}[Consideraciones]
	Al momento de realizar comparaciones, tambien se puede hacer con el valor en hexadecimal su valor en ASCII a comparar

	\begin{codebox}[ASM]{Comparacion con valor en hexadecimal}
		cmp BYTE PTR [rax], 0x70 ; 0x70 = 112 = 'p'
	\end{codebox}
\end{infobox}

\begin{flagbox}{Flag Capturada}
	pwn.college{wNjWDbSIItWSvndFqwZshHkx8O-.01N0czMywSMycjN4EzW}
\end{flagbox}

\subsection{Conditionals Without Conditionals} \label{subsec:conditionals_without_conditionals}

Para poder realizar una serie de condicionales, se suele trabajar con una \textbf{\textit{jump table}}
en vez de tener una secuencia de condicionales, una detrás de otra.

Si desamblamos un programa vamos a ver algo así:

\begin{codebox}[ASM]{}
	mov    rax, 0                    ; zero out rax
	mov    al, BYTE PTR [rcx]        ; load the character into the low byte of rax
	mov    rax, [rax*8+0x1234000]    ; load a stored address from the jump table at 0x1234000
	jmp    rax                       ; jump to it
\end{codebox}

En este caso, el valor de \textbf{0x1234000}, representa la dirección donde se encuentra la
tabla de saltos.

En este problema, nos piden:
\begin{enumerate}
	\item Encontrar la dirección de la tabla de saltos, con \mintinline{bash}{objdump}
	\item Con el GDB detectar cual es la dirección correcta a la que tiene que saltar para que salga existosamente.
	\item calcular el valor correcto, con la siguiente formula \( \text{Índice} = \frac{\text{Dirección Slot}-\text{Dirección Base}}{8} \)
\end{enumerate}
\begin{flagbox}{Flag Capturada}
	pwn.college{QuJcHz6h6ZmYjOh93E_-XKn_J3b.0FO0czMywSMycjN4EzW}
\end{flagbox}

\subsection{Recognizing Loops} \label{subsec:recognizing_loops}

\subsection{Writing Loops} \label{subsec:writing_loops}

\subsection{Writing From a Shared Library} \label{subsec:writing_from_a_shared_library}

\subsection{Returning a Value} \label{subsec:returning_a_value}

\subsection{Calling Through a Pointer} \label{subsec:calling_through_a_pointer}

\subsection{Calling Through a Pointer with an Argument} \label{subsec:calling_through_a_pointer_with_an_argument}

\subsection{Saving Caller - Saved Registers} \label{subsec:saving_caller_saved_registers}

\subsection{Saving Callee - Saved Registers} \label{subsec:saving_callee_saved_registers}
